\chapter{Géometrie des variétés algébriques}

\section{Dimension}

Soit \(k = \overline{k}\). Soit \(V\) une VA irréductible. On veut definir
\(\dim V \in \mathbb{N}\) telle que:
\begin{enumerate}[label = \roman*.]
    \item \(\dim \mathbb{A}^{n} = n\) 
    \item Si \(V\) est irréductible et \(f \in \mathcal{O}{(V)}\), \(f \neq 0 \)
        et \(f \not\in  \mathcal{O}{(V)}^{*}\), alors \(V{(f)}\) a toutes ses
        componenets irréductible de dimension \(\dim V - 1\).
    \item (en conséquence de \(ii.\)) \(\dim V\) est la longueur maximale d'une
        chaîne de fermés irréductibles
        \[
          \varnothing \subsetneq V_d \subsetneq V_{d-1} \subsetneq \dots
          \subsetneq V_{1} \subsetneq V_{0} = V
        \]
\end{enumerate}

La condition \(iii.\) est donc prise comme definition de dimension, si l'espace
est Noetherien.

\begin{example}[de ii.]
    Soit \(V = \mathbb{A}^{n}\) et \(f \in \mathcal{O}{(V)}\). Alors \(f =
    f_{1}^{a_{1}} \dots f_r^{a_r}\) une décomposition en produits d'éléments
    premiers. Alors \(V{(f)} = \bigcup_{i=1} ^{n} V{(f_{i})}\) qui ont dimension
    \(n-1\).
\end{example}

On rappel que si \(K / k \) est une extension de corps de type fini (i.e. \(K =
k{(y_{1}, \dots, y_{n})}\)), alors il existe des éléments \(x_{1}, \dots, x_d
\in K\) tels que \( k \subseteq k{(x_{1}, \dots, x_d)} \subseteq K  \) où la
première extension est transcenante pure (i.e. les \(x_{i}\) sont algébriquement
indépendents) et la deuxième est une extension finie (séparable?) de corps.
Alors \(d =: \mathrm{deg}\,\mathrm{tr}_k K\) est le dégree de trancendance de
\(K\).

\begin{definition}{(algébrique)}
    Soit \(V\) une variété algébrique sur \(k\) de type fini. Si \(V\) est
    irréductible, alors \(\dim V = \mathrm{deg}\,\mathrm{tr}_k {(\mathcal{M}{(V)})}\) 

    En général, on pose \(\dim V := \max \{\dim W, W \subseteq V \text{
    component irréductible} \} \) 
\end{definition}

\begin{example}{}
    En effet
    \[
      \dim \mathbb{A}^{n} = \mathrm{deg}\,\mathrm{tr}_k \mathrm{Frac}{\left(
      k[X_{1}, \dots, X_{n}] \right)} = k{(X_{1}, \dots, X_{n})} = n
    \]
\end{example}

